\documentclass[11pt]{article}
\newcommand{\shorttitle}{5.1: Sequence models}
\newcommand{\notesdate}{24 October 2026}
\usepackage[T1]{fontenc}
\usepackage[utf8]{inputenc}
\usepackage{lmodern}
\usepackage[a4paper,margin=23mm,headheight=15pt]{geometry}
\usepackage{amsmath,amssymb,bm,graphicx,booktabs,longtable,array}
\usepackage{xcolor,microtype,enumitem,fancyhdr,fancyvrb,titlesec,tikz}
\usepackage[breakable]{tcolorbox}
\usepackage{hyperref}
\definecolor{navy}{HTML}{163653}
\definecolor{teal}{HTML}{007F86}
\definecolor{pale}{HTML}{EEF5F7}
\hypersetup{colorlinks=true,linkcolor=navy,urlcolor=teal,pdfauthor={Study notes based on the supplied teacher slides}}
\pagestyle{fancy}\fancyhf{}
\fancyhead[L]{\small\sffamily\color{navy}Artificial Intelligence}
\fancyhead[R]{\small\sffamily\color{navy}\shorttitle}
\fancyfoot[C]{\small\thepage}
\renewcommand{\headrulewidth}{0.3pt}
\titleformat{\section}{\Large\sffamily\bfseries\color{navy}}{\thesection}{0.7em}{}
\titleformat{\subsection}{\large\sffamily\bfseries\color{teal}}{\thesubsection}{0.6em}{}
\titleformat{\subsubsection}{\normalsize\sffamily\bfseries\color{navy}}{\thesubsubsection}{0.6em}{}
\setcounter{tocdepth}{2}\setcounter{secnumdepth}{2}
\setlength{\parindent}{0pt}\setlength{\parskip}{6pt plus 1pt}
\linespread{1.06}\setlist{itemsep=3pt,topsep=4pt,leftmargin=1.5em}
\setlength{\emergencystretch}{3em}
\newtcolorbox{keyidea}[1]{breakable,colback=pale,colframe=teal,boxrule=0.6pt,arc=1mm,title={#1},fonttitle=\sffamily\bfseries,coltitle=white,colbacktitle=teal}
\newtcolorbox{careful}[1]{breakable,colback=white,colframe=navy,boxrule=0.6pt,arc=1mm,title={#1},fonttitle=\sffamily\bfseries,coltitle=navy,colbacktitle=pale}
\newcommand{\sources}[1]{{\small\color{teal}\textit{Teacher slides: #1}}\par\nopagebreak[4]}
\newcommand{\newsection}[1]{\par\begingroup\skip0=\pagegoal\advance\skip0 by-\pagetotal\ifdim\skip0<12\baselineskip\newpage\fi\endgroup\section{#1}}
\newcolumntype{P}[1]{>{\raggedright\arraybackslash}p{#1}}
\newcommand{\E}{\operatorname{E}}
\newcommand{\Var}{\operatorname{Var}}
\newcommand{\Cov}{\operatorname{Cov}}
\newcommand{\R}{\mathbb{R}}
\newcommand{\argmax}{\operatorname*{arg\,max}}
\newcommand{\argmin}{\operatorname*{arg\,min}}
\newcommand{\relu}{\operatorname{ReLU}}
\newcommand{\codeinline}[1]{\texttt{\detokenize{#1}}}
\providecommand{\notesdate}{2 October 2026}
\DefineVerbatimEnvironment{code}{Verbatim}{fontsize=\small,samepage=true,frame=single,rulecolor=\color{teal},framesep=2.5mm}
\newcommand{\cover}[3]{%
\begin{titlepage}\vspace*{12mm}
{\sffamily\color{teal}\large ARTIFICIAL INTELLIGENCE\par}\vspace{9mm}
{\sffamily\bfseries\color{navy}\Huge TOPIC #1\par}\vspace{5mm}
{\sffamily\color{navy}\LARGE #2\par}\vspace{11mm}
{\large Explanatory study notes\par}
{\large Theory, worked examples and revision questions\par}\vspace{12mm}
\begin{keyidea}{What you should be able to do}#3\end{keyidea}
\vfill
Based on the teacher PDFs supplied in \texttt{Lecture Slides}. The notes explain and reorganize that material. Worked calculations are study examples derived from the slide concepts, unless explicitly identified as a slide exercise. Clarifications of ambiguous or incorrect expressions are marked in context. References use \textbf{PDF page numbers}, counted from 1, rather than printed slide numbers.
\par\vspace{4mm}{\small English edition\enspace\textbullet\enspace \notesdate}
\end{titlepage}\setcounter{page}{2}}

\hypersetup{pdftitle={5.1 - Sequence Models and LSTM Extra Material}}
\begin{document}
\cover{5.1}{Sequence models and recurrent networks}{Distinguish learned weights from changing states; count RNN and LSTM parameters; explain backpropagation through time and gates; perform a small state update; compare recurrence, temporal convolutions and attention; explain teacher forcing.}
\tableofcontents\clearpage

\newsection{Sequential data: order and context}
\sources{5.1. Sequence Models, pp. 3--6.}
A sequence is an ordered collection $(x_1,\ldots,x_T)$. Reordering words can change a sentence's meaning; reversing video frames can change the action. A useful model must account for these relationships rather than treating each item as an unrelated sample.

Different tasks have different input/output arrangements. Sentiment classification maps many words to one label. Video labeling may map one sequence to a label at every time step. Translation maps an input sentence to an output sentence of a possibly different length. Image captioning maps an image representation to a sequence of words.

A recurrent neural network (RNN) processes one element at a time while maintaining a hidden state. The same transition is reused for each step, so the model can process different sequence lengths without introducing a new set of weights for each position.
\begin{careful}{Variable length is not unique to RNNs}
Dense layers require a specified feature dimension, but convolutional architectures can also accept variable spatial or temporal sizes when their readout permits it. The important RNN distinction here is its recurrent state and shared transition, not that every other architecture must have a fixed number of inputs.
\end{careful}

\newsection{A vanilla RNN: equations, states and predictions}
\sources{5.1. Sequence Models, pp. 6--12.}
Using column vectors, input size $D$, hidden size $H$, and output size $O$:
\[
h_t=\tanh(W_{xh}x_t+W_{hh}h_{t-1}+b_h),\qquad
s_t=W_{hy}h_t+b_y.
\]
The hidden state $h_t\in\R^H$ summarizes the current input and previous state. For classification or next-token prediction, the scores $s_t$ usually pass through softmax to obtain probabilities. Scores themselves need not lie in $[0,1]$ or sum to one.

\begin{longtable}{@{}P{.20\textwidth}P{.20\textwidth}P{.50\textwidth}@{}}
\toprule Object & Shape & Role\\\midrule\endhead
$x_t$ & $D\times1$ & Current observation or token representation.\\
$W_{xh}$ & $H\times D$ & Maps input features into the state update.\\
$W_{hh}$ & $H\times H$ & Maps the previous hidden state into the update.\\
$b_h$ & $H\times1$ & Hidden bias.\\
$h_t$ & $H\times1$ & Computed state, different at different steps.\\
$W_{hy}$ & $O\times H$ & Maps the state to output scores.\\
$b_y$ & $O\times1$ & Output bias.\\\bottomrule
\end{longtable}
\begin{keyidea}{Parameters are learned; states are computed}
During an ordinary forward pass with fixed model weights, $x_t$, $h_t$, scores and predictions can change at every step. The shared weight matrices do not change merely because the next token arrives. An optimizer updates parameters during training, after gradients have been computed. A hidden-state entry is a learned feature, not automatically a particular word or class.
\end{keyidea}

\subsection{Character-level example}
For a vocabulary of $V$ characters, a one-hot input has $D=V$ and a next-character output has $O=V$. Multiplying $W_{xh}$ by a one-hot vector selects one column of $W_{xh}$. That column is the input contribution associated with the character; the final state also depends on the preceding context through $W_{hh}h_{t-1}$.

The slide examples show units responding to patterns such as URLs or bracket structure. This illustrates that learned states can track useful context. It does not imply that every unit has a simple human-readable meaning or that the network contains a separate unit for every possible concept.

\subsection{Parameter count}
With one hidden bias and one output bias,
\[
N_{\mathrm{RNN}}=HD+H^2+H+OH+O.
\]
For $V=3,H=2$, a next-character model has $6+4+2+6+3=21$ parameters. Processing 10 rather than 100 characters changes the number of operations and stored training activations, not this count. The slide's displayed vanilla-RNN equations omit biases; if a question omits them explicitly, remove the corresponding terms.

\newsection{Training through time}
\sources{5.1. Sequence Models, pp. 13--16, 32.}
\subsection{Unrolling and shared gradients}
Unroll the recurrence over $T$ steps. This gives a computation graph with $T$ copies of the transition operation, all using the same parameters. A common objective is
\[
L=\sum_{t=1}^{T}\ell_t,\qquad
\ell_t=-\log p_\theta(y_t^*\mid x_{\leq t}).
\]
Backpropagation through time (BPTT) applies the chain rule to that unrolled graph. Because the same weight contributes at many time steps, its gradient accumulates contributions from those uses. It is still one trainable parameter.

\subsection{Why gradients vanish or explode}
A distant state's influence contains a product of Jacobians:
\[
\frac{\partial h_t}{\partial h_k}
=J_tJ_{t-1}\cdots J_{k+1},\qquad
J_t=\operatorname{diag}(1-h_t^2)\,W_{hh}.
\]
Repeated multiplication by factors with small norms can make the gradient tiny. Factors with large amplification can cause it to explode. Saturated $\tanh$ units also have small derivatives. For a scalar illustration, ten factors of $0.5$ yield $0.5^{10}\approx0.00098$; ten factors of $2$ yield $1024$.

The slides' singular-value argument is a bound: a largest singular value below one gives a contraction bound for this vanilla transition. A value above one permits growth; it does not guarantee that every gradient grows, because activation derivatives and directions also matter.

\subsection{Clipping, truncation and architectural changes}
Norm clipping rescales a gradient vector $g$ when $\|g\|>\tau$:
\[
g_{\mathrm{clip}}=g\min\left(1,\frac{\tau}{\|g\|}\right).
\]
For $g=(3,4)$ and threshold $\tau=2$, the norm is 5 and the clipped vector is $(1.2,1.6)$. Its direction is unchanged. This limits an exploding gradient; it does not recover a gradient that has vanished. The parameter step also depends on the optimizer.

Truncated BPTT limits differentiation to a window of $K$ steps. A state may be carried forward while its earlier computation history is detached. This saves memory and computation, but cuts the gradients needed to assign credit farther back. LSTM/GRU gates, suitable initialization and normalization address different aspects of recurrent optimization; they are not interchangeable with truncation.

\newsection{LSTM: a memory state controlled by gates}
\sources{5.1. Sequence Models, pp. 17--21; 5.1. Extra Material, pp. 2--4.}
An LSTM has both a cell state $c_t$ and a hidden state $h_t$, each of dimension $H$. For each gate/candidate $r\in\{f,i,g,o\}$, form
\[
a_r=W_{xr}x_t+W_{hr}h_{t-1}+b_r.
\]
Then
\[
\begin{aligned}
f_t&=\sigma(a_f), & i_t&=\sigma(a_i),\\
g_t&=\tanh(a_g), & o_t&=\sigma(a_o),\\
c_t&=f_t\odot c_{t-1}+i_t\odot g_t,&
h_t&=o_t\odot\tanh(c_t).
\end{aligned}
\]
Here $\odot$ is elementwise multiplication. The sigmoid gates lie between 0 and 1; the candidate $g_t$ lies between $-1$ and 1. The cell state is not restricted to either interval.

\begin{longtable}{@{}P{.21\textwidth}P{.69\textwidth}@{}}
\toprule Quantity & Interpretation\\\midrule\endhead
Forget gate $f_t$ & Fraction of the old cell content retained. Near 1 keeps it; near 0 removes it.\\
Input gate $i_t$ & How much of the new candidate is written.\\
Candidate $g_t$ & Proposed signed content, not itself a probability.\\
Output gate $o_t$ & How much of the transformed cell state is exposed as $h_t$.\\\bottomrule
\end{longtable}
The additive cell update supplies a path that can preserve information. Along the direct cell-to-cell path, the derivative is the forget gate: if it stays near one, that path avoids repeated multiplication by a dense recurrent matrix. Gates help with long dependencies, but do not guarantee unlimited memory or perfect gradients.

\subsection{One small update}
Suppose one unit has $c_{t-1}=2$, $f_t=0.5$, $i_t=0.25$, $g_t=0.8$, and $o_t=1$. Then
\[
c_t=0.5(2)+0.25(0.8)=1.2,\qquad h_t=\tanh(1.2).
\]
The old contribution is 1 and the newly written contribution is 0.2. Closing the output gate would hide the cell from the output without necessarily erasing the stored cell value.

\newsection{LSTM sizes and parameters: the extra-material exercise}
\sources{5.1. Extra Material, pp. 5--6.}
The class exercise uses a one-hot vocabulary $\{a,b,c\}$, so $V=3$, with $H=2$. Adopt the column-vector convention above and one combined bias vector per gate.
\begin{longtable}{@{}P{.25\textwidth}P{.20\textwidth}P{.45\textwidth}@{}}
\toprule Object & Shape & Meaning\\\midrule\endhead
$x_t$ & $3\times1$ & Selects the current letter.\\
$W_{xf},W_{hf},b_f$ & $2\times3$, $2\times2$, $2$ & Forget-gate input weights, recurrent weights and bias.\\
Stacked $W_x$ & $8\times3$ & Four blocks, each with two units.\\
Stacked $W_h$ & $8\times2$ & Four recurrent blocks.\\
Stacked bias & $8$ & Two biases in each of four blocks.\\
$c_t,h_t$ & $2\times1$ each & Computed memory and exposed state.\\
$W_y,b_y$ & $3\times2$, $3$ & Next-character output head.\\\bottomrule
\end{longtable}
The LSTM cell has
\[
\boxed{N_{\mathrm{LSTM}}=4(HD+H^2+H)=4H(D+H+1).}
\]
For $D=3,H=2$, this is $4(6+4+2)=48$. The output head adds $3(2+1)=9$, giving 57 in total. Implementations with separate input and recurrent bias vectors count extra stored bias parameters; a counting question must state its convention.

The input-specific contribution of letter $b$ is stored in column $b$ of stacked $W_x$: eight numbers here. Its contextual effect also depends on recurrent weights and the current state. Its output score has its own row in $W_y$ and output bias; input encoding and output prediction are different roles.

Adding a fourth letter while keeping $H=2$ expands $x_t$, $W_x$ and the output head. It does not change the lengths of $c_t,h_t$ or the shape of $W_h$. The cell now has $4(8+4+2)=56$ parameters. Conversely, keeping $V=3$ and doubling $H$ to 4 gives $4(12+16+4)=128$ cell parameters, plus 15 output-head parameters. Doubling hidden width is more than a doubling of the recurrent parameter count because of the $H^2$ term.

\newsection{Worked slide exercise: reading the letter b}
\sources{5.1. Extra Material, p. 7.}
This is the teacher's numerical exercise, using the supplied rounded sigmoid and $\tanh$ values. Start with $c_{t-1}=(1,0.5)$, $h_{t-1}=(0.5,0)$ and $x_t=(0,1,0)$. The one-hot input selects the second column of each input matrix.

Applying the displayed weights gives:
\[
\begin{aligned}
a_f&=(2+4(0.5),\;-4+2(0))=(4,-4),\\
a_i&=(-4+0,\;4+0)=(-4,4),\\
a_g&=(0,\;2-2(0.5))=(0,1),\\
a_o&=(0+3(0),\;0+2(0.5)-1)=(0,0).
\end{aligned}
\]
The gates and candidate are therefore approximately
\[
f=(0.98,0.02),\quad i=(0.02,0.98),\quad
g=(0,0.76),\quad o=(0.5,0.5).
\]
Update the cell elementwise:
\[
c_t=(0.98(1)+0.02(0),\;0.02(0.5)+0.98(0.76))
     =(0.98,0.7548).
\]
Using the slide's approximate $\tanh(0.98)=0.75$ and rounding $0.7548$ to about $0.76$, with $\tanh(0.76)\approx0.64$,
\[
h_t\approx(0.375,0.32).
\]
The first unit mostly keeps its old cell content and writes nothing; the second mostly erases its old content and writes the new positive candidate. Both expose about half the transformed cell content.

For the bonus head
$W_y=\left[\begin{smallmatrix}1&0\\0&1\\-1&-1\end{smallmatrix}\right]$
and zero bias, scores are approximately $(0.375,0.32,-0.695)$. Greedy prediction chooses $a$. These numbers are logits, not probabilities; softmax would preserve their ordering. Exact nonlinearities give slightly different decimals because the supplied lookup values are rounded.

\newsection{GRUs, stacked networks and bidirectional processing}
\sources{5.1. Sequence Models, pp. 21--24.}
A gated recurrent unit (GRU) uses update and reset gates, with a candidate state, but no separate LSTM-style cell state. It interpolates between an old state and a proposed state. Different notations swap the roles of $z$ and $1-z$, so interpret the stated update equation before saying what a gate value near one means.

With three affine blocks and one bias per block, the cell count is $3H(D+H+1)$, compared with $4H(D+H+1)$ for the LSTM convention above. Fewer parameters does not guarantee better accuracy; it is an architectural trade-off.

Stacking recurrent layers adds depth: the lower layer's output at time $t$ becomes an input to a higher layer. Each layer normally has its own weights. A bidirectional RNN processes a sequence in both directions and combines the resulting features. This gives access to future context within the observed input, but it cannot be used unchanged for online prediction when those future observations are unavailable.

For video, a CNN can encode individual frames and a recurrent model combine their temporal features. For captioning, an image or region representation conditions a word decoder. The task determines whether one final state, every state, or an output sequence is required.

\newsection{Temporal convolutions}
\sources{5.1. Sequence Models, pp. 25--27.}
A temporal convolutional network (TCN) applies convolutions along time. A causal layer at position $t$ uses only positions at or before $t$; the target alignment determines whether $x_t$ itself is already known. A centered filter using future observations would leak information in a forecasting task.

Dilation increases the temporal span without increasing the number of filter coefficients. For stride-one layers of kernel width $K$ and dilations $d_1,\ldots,d_L$, a simple sequential stack has receptive field
\[
R=1+(K-1)\sum_{\ell=1}^{L}d_\ell.
\]
Two layers with $K=3$, dilations 1 and 2, cover $1+2(1+2)=7$ time positions. WaveNet illustrates stacked dilated causal convolutions. Training positions can be processed in parallel when all inputs are known, but autoregressive generation still depends on previously generated samples.

\newsection{Encoder--decoder models and teacher forcing}
\sources{5.1. Sequence Models, pp. 28--31.}
An encoder reads the input sequence and produces a context representation. A decoder generates the target sequence, usually beginning with a start token and stopping at an end token. Input and output lengths need not be equal.

In basic recurrent seq2seq, the encoder's final state is a fixed-size context vector. This can become a bottleneck: information from a long input must be compressed into that one vector.

During \textbf{teacher forcing}, the decoder receives the true previous target token while learning to predict the next one:
\[
L=-\sum_t\log p_\theta(y_t^*\mid y_{<t}^*,x).
\]
At generation time, the true target is unavailable, so generated tokens are fed back. An early error can create contexts less familiar from training: exposure bias. Scheduled sampling mixes in model-generated inputs during training; it is an attempted remedy, not a guarantee.

Teacher forcing does not remove recurrence from an RNN decoder: $h_t$ still depends on $h_{t-1}$. A nonrecurrent transformer decoder can instead compute training positions in parallel under a causal mask.

\subsection{Attention reduces the fixed-context bottleneck}
Keep all encoder states and form a different context for each output step:
\[
c_t=\sum_{i=1}^{T}\alpha_{t,i}h_i,\qquad
\alpha_{t,i}\geq0,\quad \sum_i\alpha_{t,i}=1.
\]
The weights depend on the current decoding query and the encoded inputs. If two scalar values are 2 and 6, weights $(0.75,0.25)$ produce context 3. Different output tokens can attend to different source positions. Attention can be added to an RNN; a transformer goes further by building its main sequence processing from attention blocks.

\newsection{Revision questions and answers}
\subsection{Questions}
\begin{enumerate}
\item Which changes during a forward pass: the recurrent weights or the hidden state?
\item How many parameters has an RNN with $D=3,H=2,O=3$, including both biases?
\item Why does an unrolled RNN have more computation but not more parameters when $T$ increases?
\item What does BPTT do with gradient contributions from repeated uses of a shared weight?
\item Does clipping fix vanishing gradients? Does truncated BPTT preserve arbitrarily long credit assignment?
\item What do $f_t\approx1$ and $i_t\approx0$ do to the LSTM cell?
\item Compute $c_t$ for $c_{t-1}=2,f_t=0.5,i_t=0.25,g_t=0.8$.
\item How many cell parameters has an LSTM with $D=3,H=2$ and one bias per gate?
\item Does adding one vocabulary symbol necessarily increase the hidden-state dimension?
\item Why is a bidirectional RNN unsuitable for strict online forecasting?
\item What does teacher forcing supply to the decoder? What changes at inference?
\item What does attention change in a basic fixed-context seq2seq model?
\end{enumerate}
\subsection{Answers}
\begin{enumerate}
\item The hidden state; weights are shared and fixed within the ordinary forward pass.
\item $6+4+2+6+3=21$.
\item More applications of the same transition are computed.
\item They are accumulated into the gradient of that one shared parameter.
\item No to both: clipping caps large gradients; truncation cuts long gradient paths.
\item Mostly preserves the old cell and writes little new content.
\item $1+0.2=1.2$.
\item $4(6+4+2)=48$, excluding an output head.
\item No; input/output vocabulary dimensions change, while hidden width can stay fixed.
\item Its reverse-direction state depends on future observations.
\item The true previous token during training; generated tokens during inference.
\item Each output can access a weighted combination of all encoder states instead of relying solely on one fixed context.
\end{enumerate}

\newsection{Source guide}
\texttt{5.1. Sequence Models.pdf}: sequence structure and vanilla RNNs (pp. 3--12), BPTT and gradient problems (pp. 13--16), LSTM/GRU and variants (pp. 17--24), TCNs (pp. 25--27), seq2seq and attention (pp. 28--32), and the character-model exercise (p. 33).

\texttt{5.1. Extra Material.pdf}: cell diagrams (pp. 2--4), sizes/counts (p. 5), parameters versus states (p. 6), and the worked numerical update (p. 7). These notes incorporate this extra deck into 5.1 rather than treating it as a different course topic.
\end{document}
